\subsection{SPARQL Query Demonstration}
\label{subsec:sparql-demo}

To validate the queryability of the knowledge graph, four SPARQL queries were executed against the RDF representation. Each query targets a distinct retrieval pattern relevant to FRIA evidence gathering under Article~27 of the AI Act. Table~\ref{tab:sparql-summary} summarises the results.

\begin{table}[htbp]
\centering
\caption{SPARQL query demonstration results}
\label{tab:sparql-summary}
\begin{tabular}{clr}
\toprule
Query & Retrieval pattern & Results \\
\midrule
Q1 & Essential-services incidents involving non-discriminati & 60 \\
Q2 & Incidents where hybrid method identified profiling/scor & 47 \\
Q3 & Rights implicated by employment-domain incidents, ranke & 4 \\
Q4 & Incidents involving both privacy/data-protection and un & 12 \\
\bottomrule
\end{tabular}
\end{table}

\paragraph{Essential-services incidents involving non-discrimination rights}

\begin{lstlisting}[language=SPARQL,basicstyle=\ttfamily\footnotesize,breaklines=true,caption={Essential-services incidents involving non-discrimination rights}]
PREFIX fria: <https://example.org/fria-kg/>
PREFIX rdfs: <http://www.w3.org/2000/01/rdf-schema#>

SELECT ?incident ?label
WHERE {
    ?incident a fria:Incident ;
              rdfs:label ?label ;
              fria:hasDomain fria:domain_essential_services ;
              fria:hasRight fria:right_non_discrimination .
}
ORDER BY ?label
\end{lstlisting}

Query~1 demonstrates combined filtering across domain and rights axes. The 60 results confirm that the graph structure supports the retrieval pattern most directly required by Article~27(1)(a): identifying precedent incidents within a specific high-risk domain that implicate a particular fundamental right.

\paragraph{Incidents where hybrid method identified profiling/scoring}

\begin{lstlisting}[language=SPARQL,basicstyle=\ttfamily\footnotesize,breaklines=true,caption={Incidents where hybrid method identified profiling/scoring}]
PREFIX fria: <https://example.org/fria-kg/>
PREFIX rdfs: <http://www.w3.org/2000/01/rdf-schema#>

SELECT ?incident ?label ?domain
WHERE {
    ?incident a fria:Incident ;
              rdfs:label ?label ;
              fria:hasDomain ?domainNode ;
              fria:hasPattern fria:pattern_profiling_scoring .
    BIND(STRAFTER(STR(?domainNode), "fria-kg/") AS ?domain)
}
ORDER BY ?domain ?label
\end{lstlisting}

Query~2 retrieves incidents by system pattern, returning 47 records classified as profiling or scoring systems. This pattern-based retrieval enables deployers to locate precedents involving similar AI architectures, supporting the risk-by-analogy reasoning that underpins impact assessment.

\paragraph{Rights implicated by employment-domain incidents, ranked by frequency}

\begin{lstlisting}[language=SPARQL,basicstyle=\ttfamily\footnotesize,breaklines=true,caption={Rights implicated by employment-domain incidents, ranked by frequency}]
PREFIX fria: <https://example.org/fria-kg/>
PREFIX rdfs: <http://www.w3.org/2000/01/rdf-schema#>

SELECT ?rightLabel (COUNT(?incident) AS ?count)
WHERE {
    ?incident a fria:Incident ;
              fria:hasDomain fria:domain_employment ;
              fria:hasRight ?right .
    BIND(STRAFTER(STR(?right), "fria-kg/") AS ?rightLabel)
}
GROUP BY ?rightLabel
ORDER BY DESC(?count)
\end{lstlisting}

Query~3 aggregates rights across employment-domain incidents, demonstrating cross-axis analytical capability. The ranked output reveals which fundamental rights are most frequently implicated in employment AI systems, providing quantitative evidence for FRIA prioritisation decisions.

\paragraph{Incidents involving both privacy/data-protection and unfair-exclusion harms}

\begin{lstlisting}[language=SPARQL,basicstyle=\ttfamily\footnotesize,breaklines=true,caption={Incidents involving both privacy/data-protection and unfair-exclusion harms}]
PREFIX fria: <https://example.org/fria-kg/>
PREFIX rdfs: <http://www.w3.org/2000/01/rdf-schema#>

SELECT ?incident ?label ?domain
WHERE {
    ?incident a fria:Incident ;
              rdfs:label ?label ;
              fria:hasDomain ?domainNode ;
              fria:hasHarm fria:harm_privacy_breach ;
              fria:hasHarm fria:harm_unfair_exclusion .
    BIND(STRAFTER(STR(?domainNode), "fria-kg/") AS ?domain)
}
ORDER BY ?label
\end{lstlisting}

Query~4 identifies incidents exhibiting both privacy breach and unfair exclusion harms simultaneously, returning 12 records. This multi-label intersection query demonstrates that the graph's multi-valued harm annotations support compound risk identification --- a capability not available in flat tabular representations.
